\chapter{Encapsulamiento y Composición}

\section{Encapsulamiento y Ocultamiento de Información}
El \textit{encapsulamiento} es el principio de la POO que consiste en ocultar los detalles internos de implementación
de un objeto y restringir el acceso directo a sus atributos desde el exterior del código de la clase.
Esto evita que los datos internos sean modificados de forma arbitraria o inconsistente.

A diferencia de lenguajes como C++ o Java, Python no posee palabras clave como \texttt{private} o \texttt{protected}
a nivel de compilador, sino que emplea convenciones de nombres para indicar la privacidad de los miembros.

\subsection{Atributos Protegidos (\texttt{\_atributo})}
Un atributo cuyo nombre comienza con un único guion bajo (por ejemplo, \texttt{\_balance}) indica convencionalmente
que es un miembro protegido. Aunque Python permite acceder a él desde el exterior, la regla indica que solo debería
ser utilizado internamente por la clase y sus subclases.

\subsection{Atributos Privados (\texttt{\_\_atributo}) y Name Mangling}
Cuando un atributo comienza con dos guiones bajos (por ejemplo, \texttt{\_\_balance}), Python activa un mecanismo
llamado \textit{Name Mangling} (modificación de nombres). El interprete transforma internamente el nombre del atributo
agregando el prefijo \texttt{\_NombreDeClase}, lo que dificulta su acceso o modificación accidental desde afuera.

\begin{lstlisting}[caption={Diferencia entre atributos protegidos y privados.}, label={list:name_mangling}]
class BankAccount:
    def __init__(self, owner, balance):
        self.owner = owner
        self._protected_data = "Dato protegido"
        self.__balance = balance  # Atributo privado

cuenta = BankAccount("Laura", 5000)

print(cuenta.owner)            # Acceso directo permitido
print(cuenta._protected_data)  # Permitido por sintaxis, pero desaconsejado

# El acceso a cuenta.__balance arroja AttributeError
# Para acceder forzadamente (mangling): print(cuenta._BankAccount__balance)
\end{lstlisting}

\section{Propiedades (\texttt{@property} y Setters)}
En lugar de escribir métodos de acceso tradicionales al estilo Java (como \texttt{get\_balance()} y \texttt{set\_balance()}),
Python ofrece una alternativa elegante y concisa mediante decoradores de \textit{propiedades}: \texttt{@property} para el getter
y \texttt{@propiedad.setter} para el setter.

Las propiedades permiten acceder y modificar atributos privados mediante la notación habitual de atributos (\texttt{obj.balance = 100}),
pero ejecutando código de validación transparente detrás de escena.

\begin{lstlisting}[caption={Uso de \texttt{@property} para validar la asignación de saldo.}, label={list:property_setter}]
class BankAccount:
    def __init__(self, balance):
        self._balance = balance

    @property
    def balance(self):
        """Getter: retorna el saldo protegido."""
        return self._balance

    @balance.setter
    def balance(self, amount):
        """Setter: valida que el saldo no sea negativo."""
        if amount < 0:
            raise ValueError("El saldo no puede ser negativo.")
        self._balance = amount

cuenta = BankAccount(1000)
print(f"Saldo inicial: {cuenta.balance}")  # Invoca el getter

cuenta.balance = 1500                     # Invoca el setter con validacion
print(f"Nuevo saldo: {cuenta.balance}")
\end{lstlisting}

\section{Relaciones entre Objetos: Composición y Agregación}
Mientras que la herencia modela una relación del tipo \textit{``es-un''} (un colectivo \textit{es un} vehículo),
la \textit{composición} y la \textit{agregación} modelan relaciones del tipo \textit{``tiene-un''}.

La composición ocurre cuando una clase compleja contiene instancias de otras clases como atributos propios,
coordinando la interacción entre dichos objetos.

\begin{figure}[!ht]
  \centering
  \begin{tikzpicture}[node distance=2cm, align=center]
    \node [draw, rectangle, fill=blue!20, minimum width=3.5cm, minimum height=1.5cm] (vuelo) 
      {\textbf{Contenedor}\\ \texttt{Flight}};
    
    \node [draw, rectangle, fill=green!20, minimum width=3.5cm, minimum height=1.2cm, right=2.5cm of vuelo] (pasajero) 
      {\textbf{Componente}\\ \texttt{Passenger}};
    
    \draw[->, thick, double] (vuelo) -- node[above] {\small Contiene lista de} (pasajero);
  \end{tikzpicture}
  \caption{Relación de Composición: Un Vuelo contiene una lista de Pasajeros.}
  \label{fig:composicion_vuelo}
\end{figure}

A continuación se presenta una clase \texttt{Flight} que gestiona la reserva y capacidad ocupada por objetos \texttt{Passenger}:

\begin{lstlisting}[caption={Ejemplo de Composición con clases Flight y Passenger.}, label={list:composicion_flight}]
class Passenger:
    def __init__(self, name, passport_number):
        self.name = name
        self.passport_number = passport_number

class Flight:
    def __init__(self, flight_number, capacity):
        self.flight_number = flight_number
        self.capacity = capacity
        self.passengers = []  # Lista que almacenara objetos Passenger

    def add_passenger(self, passenger):
        if len(self.passengers) < self.capacity:
            self.passengers.append(passenger)
            print(f"Pasajero {passenger.name} abordado con exito.")
            return True
        else:
            print("Vuelo completo. No se pudo agregar al pasajero.")
            return False

p1 = Passenger("Carlos Gomez", "AB123456")
p2 = Passenger("Maria Lopez", "CD789012")

vuelo = Flight("AR1302", 2)
vuelo.add_passenger(p1)
vuelo.add_passenger(p2)
\end{lstlisting}

Del mismo modo, la composición permite implementar patrones donde una clase administradora procesa de forma polimórfica
una colección de objetos contenidos en su interior:

\begin{lstlisting}[caption={Clase Zoo administrando una colección polimórfica de animales.}, label={list:zoo_animales}]
class Animal:
    def __init__(self, name):
        self.name = name

    def eat(self):
        return f"{self.name} esta comiendo."

class Zoo:
    def __init__(self):
        self.animals = []

    def add_animal(self, animal):
        self.animals.append(animal)

    def feed_all(self):
        for animal in self.animals:
            print(animal.eat())

zoo = Zoo()
zoo.add_animal(Animal("Leon"))
zoo.add_animal(Animal("Panda"))
zoo.feed_all()
\end{lstlisting}
